\section{Blade element corrections and the figure of merit}
\label{sec:blade}

Momentum theory bounds the power from below. A real rotor pays two further
tolls: the induced flow is not uniform and the blades are finite in number, and
the blade sections have profile drag. This section derives the second from
blade element theory and absorbs the first into a coefficient, with the
consequence that the figure of merit becomes an output of blade geometry rather
than an input.

\subsection{Induced power factor}

\begin{model}[Non-ideal induced power]
\label{mod:kappa}
The induced power is written
\begin{equation}
  \Pind = \kappa\,T\vind, \qquad \kappa > 1 ,
  \label{eq:kappa}
\end{equation}
where $\kappa$ collects non-uniform inflow, tip loss and swirl. Typical values
for small rotors are $\kappa \in [1.10,1.25]$~\cite[\S2.4]{leishman2006}; we
take $\kappa = 1.15$.
\end{model}

The alternative is Prandtl's tip-loss factor $B$, with the induced velocity
computed over an effective disk of radius $BR$. The two are related by
$\kappa \approx B^{-2}$ for small losses; \cref{eq:kappa} is used because it
keeps the algebra of \cref{sec:closure} in closed form.

\subsection{Profile power}

\begin{definition}[Solidity]
For a rotor with $N_b$ blades of constant chord $c$ and radius $R$, the
solidity is the ratio of blade planform area to disk area,
\begin{equation}
  \sigma \coloneqq \frac{N_b c R}{\pi R^2} = \frac{N_b c}{\pi R}
        = \frac{N_b}{\pi\,\mathrm{AR}},
  \qquad \mathrm{AR} \coloneqq \frac{R}{c} .
  \label{eq:solidity}
\end{equation}
\end{definition}

\Cref{fig:element} fixes the notation of blade element theory.

\begin{figure}[tbp]
\centering
% Blade element theory. (a) the rotor seen from above: two blades of chord c,
% the annulus at radius r of width dr and the element it cuts from a blade.
% (b) section A-A seen from the blade tip: velocity triangle, angles and the
% resolution of lift and drag into thrust and in-plane force.
\begin{tikzpicture}[
  declare function={yt(\s) = 0.6*(0.2969*sqrt(\s) - 0.1260*\s
                     - 0.3516*\s*\s + 0.2843*\s*\s*\s - 0.1015*\s*\s*\s*\s);}]
  % ------------------------------------------------------------ panel (a)
  \begin{scope}
    \draw[cG, dashed] (0,0) circle (1.9);
    \fill[cB!10, even odd rule] (0,0) circle (1.36) circle (1.14);
    \draw[cB!60, line width=0.3pt] (0,0) circle (1.36) (0,0) circle (1.14);
    \foreach \s in {1,-1}{
      \draw[cA, fill=cA!12, line width=0.7pt]
          ({\s*0.26},-0.315) rectangle ({\s*1.9},0.315);
    }
    \fill[cB!70] (1.14,-0.315) rectangle (1.36,0.315);
    \fill[black!65] (0,0) circle (0.26);
    % section marks A-A
    \draw[line width=0.9pt] (1.25,0.45) -- (1.25,0.68);
    \node[lbl, anchor=east] at (1.2,0.66) {A};
    \draw[line width=0.9pt] (1.25,-0.45) -- (1.25,-0.68);
    \node[lbl, anchor=east] at (1.2,-0.66) {A};
    % dimensions
    \draw[dim] (0,-1.05) -- node[lbl, below, fill=white, inner sep=1pt] {$r$} (1.25,-1.05);
    \draw[dim] (1.14,1.02) -- (1.36,1.02);
    \node[lbl, above] at (1.25,1.06) {$\dd r$};
    \draw[dim] (1.72,-0.315) -- node[lbl, right] {$c$} (1.72,0.315);
    \draw[dim] (0,0) -- node[lbl, above left, inner sep=1pt] {$R$} (125:1.9);
    % rotation, clockwise seen from above
    \draw[vec, cA] (80:2.2) arc (80:35:2.2) node[right, lbl] {$\Omega$};
    \node[font=\small] at (0,-2.5) {(a) rotor seen from above};
  \end{scope}
  % ------------------------------------------------------------ panel (b)
  \begin{scope}[xshift=8.3cm, yshift=-0.2cm]
    \def\ph{12} \def\th{20}
    % plane of rotation
    \draw[construct] (-4.2,0) -- (3.6,0) node[right, lbl, align=left] {plane of\\rotation};
    % aerofoil, quarter chord at the origin, leading edge left, nose up by theta
    \begin{scope}[rotate=-\th]
      \draw[cA, fill=cA!15, line width=0.8pt]
        plot[domain=0:180, samples=70] ({3.0*((1-cos(\x))/2-0.25)}, { 3.0*yt((1-cos(\x))/2)})
        -- plot[domain=180:0, samples=70] ({3.0*((1-cos(\x))/2-0.25)}, {-3.0*yt((1-cos(\x))/2)})
        -- cycle;
      \draw[cG, dash dot, line width=0.4pt] (-3.5,0) -- (2.6,0);
    \end{scope}
    \node[lbl, cG, anchor=north] at ({-3.6*cos(\th)},{3.6*sin(\th)-0.02}) {chord};
    % direction of the resultant wind through the section
    \draw[construct] ({-3.6*cos(\ph)},{3.6*sin(\ph)}) -- (0,0);
    % angles on the upstream side
    \draw[cG] (180:1.55) arc (180:{180-\ph}:1.55);
    \node[lbl] at ({180-\ph/2}:1.82) {$\phi$};
    \draw[cG] ({180-\ph}:2.3) arc ({180-\ph}:{180-\th}:2.3);
    \node[lbl] at ({180-(\ph+\th)/2}:2.58) {$\alpha$};
    \draw[cG] (180:3.05) arc (180:{180-\th}:3.05);
    \node[lbl] at ({180-\th/2+1}:3.35) {$\theta$};
    % velocity triangle, drawn apart
    \begin{scope}[shift={(-4.1,1.95)}]
      \pgfmathsetmacro{\tp}{2.2*tan(\ph)}
      \draw[vec, cB] (0,0) -- node[lbl, above] {$U_T = \Omega r$} (2.2,0);
      \draw[vec, cB] (2.2,0) -- node[lbl, right] {$U_P = \vind$} (2.2,-\tp);
      \draw[vec, black, line width=1.0pt] (0,0) -- node[lbl, below left, inner sep=1pt] {$U$} (2.2,-\tp);
      \draw[cG] (0.75,0) arc (0:-\ph:0.75);
      \node[lbl] at (-6:1.0) {$\phi$};
    \end{scope}
    % forces at the quarter chord
    \pgfmathsetmacro{\Lx}{2.4*sin(\ph)} \pgfmathsetmacro{\Ly}{2.4*cos(\ph)}
    \pgfmathsetmacro{\Dx}{0.75*cos(\ph)} \pgfmathsetmacro{\Dy}{-0.75*sin(\ph)}
    \coordinate (F) at ({\Lx+\Dx},{\Ly+\Dy});
    \draw[construct] (\Lx,\Ly) -- (F) -- (\Dx,\Dy);
    \draw[construct] (F) -- (0,{\Ly+\Dy});
    \draw[construct] (F) -- ({\Lx+\Dx},0);
    \draw[vec, cA, line width=1.0pt] (0,0) -- (\Lx,\Ly) node[above right, lbl, inner sep=1pt] {$\dd L$};
    \draw[vec, cD, line width=1.0pt] (0,0) -- (\Dx,\Dy);
    \node[lbl, cD, anchor=north west, inner sep=1pt] at (\Dx,\Dy) {$\dd D$};
    \draw[vec, black] (0,0) -- (F);
    \draw[vec, cC!70!black, line width=1.1pt] (0,0) -- (0,{\Ly+\Dy})
        node[left, lbl] {$\dd T$};
    \draw[vec, cC!70!black, line width=1.1pt] (0,0) -- ({\Lx+\Dx},0);
    \node[lbl, cC!60!black, anchor=south west, inner sep=1pt] at ({\Lx+\Dx},0.02) {$\dd Q/r$};
    % blade motion
    \draw[vec, cG] (-0.6,-1.2) -- (-2.2,-1.2) node[left, lbl] {blade motion};
    \node[font=\small] at (-0.4,-2.3) {(b) section A-A, seen from the tip};
  \end{scope}
\end{tikzpicture}

\caption{Blade element notation. (a) A two-bladed rotor of radius $R$ and chord
$c$, so $\sigma = N_bc/(\pi R)$; the annulus at radius $r$ of width $\dd r$
cuts an element from each blade. (b) The element in section. It meets the air
at $U_T = \Omega r$ in the plane of rotation and $U_P = \vind$ through it, so
the resultant $U$ is inclined at the inflow angle $\phi = \arctan(U_P/U_T)$
and the angle of attack is $\alpha = \theta - \phi$ for a pitch $\theta$. Lift
$\dd L$ is normal to $U$ and drag $\dd D$ along it. Resolved on the rotor axes
they give the thrust $\dd T = \dd L\cos\phi - \dd D\sin\phi$ and the in-plane
force $\dd Q/r = \dd L\sin\phi + \dd D\cos\phi$. \Cref{prop:profile} keeps
the drag part of the second in the limit $\phi\to0$, where
$\dd P = \Omega r\,\dd D$.}
\label{fig:element}
\end{figure}

\begin{proposition}[Profile power of a rotor in hover]
\label{prop:profile}
Let each blade section have drag coefficient $\Cdz$, taken constant along the
span, and let the rotor turn at angular rate $\Omega$ with tip speed
$\Vtip = \Omega R$. Neglecting the inflow angle in the drag direction, the
profile power is
\begin{equation}
  \Pprof = \frac{1}{8}\,\rho\, A\, \Vtip^{3}\, \sigma\, \Cdz .
  \label{eq:profile}
\end{equation}
\end{proposition}

\begin{proof}
Consider a blade element at radius $r$ of span $\dd r$. Its sectional velocity
is $\Omega r$ to leading order, so the drag on the element is
\[
  \dd D = \tfrac12 \rho (\Omega r)^2 c\, \Cdz \dd r ,
\]
and the power absorbed is $\dd P = \Omega r \dd D$. Integrating over the blade
and multiplying by $N_b$ blades,
\[
  \Pprof = N_b \int_0^R \tfrac12 \rho \Omega^3 r^3 c\,\Cdz \dd r
         = \tfrac{1}{8}\rho\,\Omega^3 c\, N_b\, \Cdz R^4 .
\]
Substituting $c = \sigma \pi R/N_b$ from \cref{eq:solidity},
\[
  \Pprof = \tfrac{1}{8}\rho\,\Omega^3 \Cdz \,\sigma \pi R^{5}
         = \tfrac{1}{8}\rho\, \Cdz\, \sigma\, (\pi R^2)(\Omega R)^3
         = \tfrac{1}{8}\rho\, A\, \Vtip^3\, \sigma\, \Cdz . \qedhere
\]
\end{proof}

\subsection{Figure of merit}

\begin{definition}[Figure of merit]
\begin{equation}
  \FM \coloneqq \frac{\Pideal}{\Pshaft}
      = \frac{T\vind}{\kappa T\vind + \Pprof} \in (0,1) .
  \label{eq:fm}
\end{equation}
\end{definition}

\begin{proposition}[Figure of merit at fixed blade loading]
\label{prop:fmblade}
Define the blade loading coefficient
\begin{equation}
  \frac{\Ct}{\sigma} \coloneqq \frac{T}{\rho A \Vtip^2 \sigma} .
  \label{eq:ctsigma}
\end{equation}
Then
\begin{equation}
  \FM = \left[\kappa + \frac{\Cdz}{8}\,
        \frac{1}{(\Ct/\sigma)}\,\frac{\Vtip}{\vind}\right]^{-1},
  \qquad
  \frac{\Vtip}{\vind} = \sqrt{\frac{2}{\sigma\,(\Ct/\sigma)}} ,
  \label{eq:fmclosed}
\end{equation}
so that at fixed blade loading the figure of merit depends on the blade
geometry only through the group $\sigma\,(\Ct/\sigma) = \Ct$, and is
independent of rotor size and of tip speed separately.
\end{proposition}

\begin{proof}
From \cref{eq:profile,eq:ctsigma}, $\rho A \sigma = T/(\Vtip^2 (\Ct/\sigma))$
and hence
\[
  \Pprof = \tfrac18 \Vtip^3 \Cdz \cdot \frac{T}{\Vtip^2 (\Ct/\sigma)}
         = \frac{\Cdz}{8}\,\frac{T\,\Vtip}{(\Ct/\sigma)} .
\]
Dividing \cref{eq:fm} through by $T\vind$ gives the first identity. For the
second, \cref{eq:vi} gives $\vind^2 = T/(2\rho A)$ and \cref{eq:ctsigma} gives
$\Vtip^2 = T/(\rho A \sigma (\Ct/\sigma))$, so
$\Vtip^2/\vind^2 = 2/(\sigma (\Ct/\sigma))$.
\end{proof}

\begin{corollary}[Design rule]
\label{cor:designrule}
Holding blade loading at its design value and increasing solidity lowers the
tip speed required to carry a given thrust,
$\Vtip \propto \sigma^{-1/2}$, and lowers the profile power with it,
$\Pprof \propto \Vtip \propto \sigma^{-1/2}$.
\end{corollary}

\begin{proof}
Immediate from \cref{eq:ctsigma} with $T$, $\rho A$ and $\Ct/\sigma$ held
fixed, together with
\[
  \Pprof = \frac{\Cdz}{8}\,\frac{T\,\Vtip}{\Ct/\sigma} . \qedhere
\]
\end{proof}

\Cref{cor:designrule} is not obvious and is used heavily later: a fatter,
slower blade is better on \emph{both} counts, quieter by
\cref{sec:acoustics} and cheaper in profile power, until the assumptions break.
They break in two places, and both are enforced as constraints.
\Cref{fig:designrule} shows both effects at the design's thrust and the two
limits that stop them.

\begin{figure}[tbp]
\centering
% Cor. design rule: at fixed blade loading Ct/sigma the figure of merit
% rises and the tip speed falls as solidity grows, eq:fmclosed and
% eq:ctsigma, at the design's hover thrust per rotor and induced factor.
\providecommand{\aThover}{4.23}
\providecommand{\aVstall}{24.71}
\providecommand{\aVdesign}{26.69}
\providecommand{\aPstall}{17.04}
\providecommand{\aPdesign}{17.56}
\providecommand{\aPind}{14.87}
\providecommand{\aCtMax}{0.14}
\providecommand{\aCtDesign}{0.12}
\providecommand{\aSigma}{0.212}
\providecommand{\aRho}{1.225}
\providecommand{\aDL}{22.2}
\providecommand{\aWake}{6.02}
\providecommand{\aKappa}{1.168}

\begin{tikzpicture}[
  declare function={
    FMc(\s,\c) = 1/(\aKappa + (\rCdzero/(8*\c))*sqrt(2/(\s*\c)));
    Ar = pi*(\rDiam/2000)^2;
    Vt(\s,\c) = sqrt(\aThover/(\aRho*Ar*\s*\c));}]
\begin{groupplot}[group style={group size=2 by 1, horizontal sep=1.55cm},
    halfplot, xmin=0.03, xmax=0.35, domain=0.03:0.35, samples=80,
    xlabel={solidity $\sigma$}]
  \nextgroupplot[ylabel={figure of merit $\FM$}, ymin=0.3, ymax=0.8,
                 legend columns=2,
                 legend style={font=\scriptsize, at={(0.02,0.03)}, anchor=south west}]
    \fill[cG!12] (axis cs:0.05,0.3) rectangle (axis cs:0.12,0.8);
    \node[lbl, cG, anchor=north] at (axis cs:0.085,0.795) {helicopters};
    \fill[cD!10] (axis cs:0.25,0.3) rectangle (axis cs:0.35,0.8);
    \node[lbl, cD, rotate=90, anchor=center] at (axis cs:0.30,0.5) {cascade, $\sigma > 0.25$};
    \addplot[cA!40] {FMc(x,0.08)}; \addlegendentry{$\Ct/\sigma = 0.08$}
    \addplot[cA!65] {FMc(x,0.10)}; \addlegendentry{$0.10$}
    \addplot[cA]    {FMc(x,0.12)}; \addlegendentry{$0.12$, design}
    \addplot[cB, dashed] {FMc(x,0.14)}; \addlegendentry{$0.14$, stall}
    \addplot[only marks, mark=*, mark size=2pt, cA, forget plot]
        coordinates {(\aSigma,{FMc(\aSigma,0.12)})};
  \nextgroupplot[ylabel={tip speed $\Vtip$ [\si{\metre\per\second}]},
                 ymin=10, ymax=95]
    \fill[cD!10] (axis cs:0.25,10) rectangle (axis cs:0.35,95);
    \addplot[cA!40] {Vt(x,0.08)};
    \addplot[cA!65] {Vt(x,0.10)};
    \addplot[cA]    {Vt(x,0.12)};
    \addplot[cB, dashed] {Vt(x,0.14)};
    \addplot[only marks, mark=*, mark size=2pt, cA]
        coordinates {(\aSigma,{Vt(\aSigma,0.12)})};
    \node[lbl, anchor=south west, fill=white, inner sep=1pt]
        at (axis cs:\aSigma,{Vt(\aSigma,0.12)+2})
        {design, \SI{\rVtip}{\metre\per\second}};
    \node[lbl, cG, anchor=north east, fill=white, inner sep=1pt]
        at (axis cs:0.34,93) {$\Vtip\propto\sigma^{-1/2}$};
\end{groupplot}
\end{tikzpicture}

\caption{The design rule of \cref{cor:designrule} at the thrust per rotor of
the design, $T = \SI{\aThover}{\newton}$, with its induced factor
$\kappa = \aKappa$ and $\Cdz = \rCdzero$ held fixed. Left: at fixed blade
loading the figure of merit \cref{eq:fmclosed} rises with solidity, because
the profile term falls as $\sigma^{-1/2}$. Right: the tip speed that carries
the thrust falls as $\sigma^{-1/2}$, \cref{eq:ctsigma}, which is what makes
the rotor quiet (\cref{sec:acoustics}). Both gains stop at the cascade limit
$\sigma = 0.25$ of \cref{as:sigma} and at the stall limit $\Ct/\sigma = 0.14$
of \cref{as:stall}. The design ($\bullet$) sits at $\sigma = \rSolidity$ and
$\Ct/\sigma = \rCtSigma$; helicopter rotors occupy the grey band.}
\label{fig:designrule}
\end{figure}

\subsection{Two limits on solidity and blade loading}

\begin{assumption}[Blade-element independence]
\label{as:sigma}
Blade element theory treats each section as an isolated aerofoil in a uniform
inflow. This fails once the blades operate as a cascade. Helicopter rotors run
$\sigma \in [0.05,0.12]$; high-solidity propellers reach about $0.25$. We
impose $\sigma \le \sigma_{\max} = 0.25$ and report violations, because beyond
it the model flatters the rotor and one no longer has a propeller but a ducted
fan without a duct.
\end{assumption}

\begin{assumption}[Stall margin]
\label{as:stall}
Sectional lift coefficient is bounded by stall, which through
$\Ct/\sigma \approx C_L/6$ for linearly twisted blades bounds the blade
loading. We impose $\Ct/\sigma \le 0.14$ and design at $0.12$.
\end{assumption}

\begin{remark}[Fixed-pitch rotors run at constant $\Ct$]
For a fixed-pitch rotor at fixed collective, $T \propto \Omega^2$ so
$\Ct = T/(\rho A \Vtip^2)$ is invariant under changes of rotor speed. The
blade-loading constraint therefore binds identically in hover and at maximum
thrust, which is why the engine reports the same $\Ct/\sigma$ at both
conditions. Stall margin cannot be recovered by spinning faster.
\end{remark}

\subsection{Reynolds number}

Both \cref{prop:profile,prop:fmblade} take $\Cdz$ constant. A design driven
toward low tip speed by \cref{cor:designrule} operates at chord Reynolds
numbers
\begin{equation}
  \mathrm{Re} = \frac{0.7\,\Vtip\, c}{\nu} ,
  \label{eq:reynolds}
\end{equation}
evaluated at the $0.7R$ station, of order $10^5$ or below, where profile drag
rises sharply. Treating $\Cdz$ as constant would hand the optimiser a free
lunch: it would drive the tip speed arbitrarily low. We therefore use
\begin{equation}
  \Cdz(\mathrm{Re}) = \Cdz^{\mathrm{ref}}
    \min\!\left\{\left(\frac{\mathrm{Re}_{\mathrm{ref}}}{\mathrm{Re}}
    \right)^{0.2},\,3\right\},
  \qquad \mathrm{Re}_{\mathrm{ref}} = 2\times 10^5 ,
  \label{eq:recorrection}
\end{equation}
an empirical power law of the standard form for attached turbulent flow, capped
to keep it from diverging. This is the crudest submodel in the paper and is
flagged as such; its only job is to prevent an unphysical optimum.
\Cref{fig:tipspeed} shows the resulting trade for one rotor of the design.

\begin{figure}[tbp]
\centering
% One rotor of the design at its hover thrust, fixed geometry, tip speed
% swept with automatic tip speed off (engine data, gen_a.py).
\providecommand{\aThover}{4.23}
\providecommand{\aVstall}{24.71}
\providecommand{\aVdesign}{26.69}
\providecommand{\aPstall}{17.04}
\providecommand{\aPdesign}{17.56}
\providecommand{\aPind}{14.87}
\providecommand{\aCtMax}{0.14}
\providecommand{\aCtDesign}{0.12}
\providecommand{\aSigma}{0.212}
\providecommand{\aRho}{1.225}
\providecommand{\aDL}{22.2}
\providecommand{\aWake}{6.02}
\providecommand{\aKappa}{1.168}

\begin{tikzpicture}
\begin{axis}[paperplot, xmin=18, xmax=80, ymin=0, ymax=80,
    xlabel={tip speed $\Vtip$ [\si{\metre\per\second}]},
    ylabel={power per rotor [\si{\watt}]},
    legend style={at={(0.3,0.97)}, anchor=north west}]
  \fill[cD!10] (axis cs:18,0) rectangle (axis cs:\aVstall,80);
  \node[lbl, cD, rotate=90, anchor=center] at (axis cs:{(18+\aVstall)/2},50)
      {stall: $\Ct/\sigma > \aCtMax$};
  \draw[cA, dashed, line width=0.6pt] (axis cs:\aVdesign,0) -- (axis cs:\aVdesign,80);
  \addplot[black] table[x=vtip, y=Pshaft] {results/fig/a_tipspeed.dat};
  \addlegendentry{shaft, $\Pind + \Pprof$}
  \addplot[cA, dashed] table[x=vtip, y=Pind] {results/fig/a_tipspeed.dat};
  \addlegendentry{induced, $\kappa T\vind$}
  \addplot[cB] table[x=vtip, y=Pprof] {results/fig/a_tipspeed.dat};
  \addlegendentry{profile, $\tfrac18\rho A\Vtip^3\sigma\Cdz(\mathrm{Re})$}
  \addplot[cG, dotted, line width=1.2pt] table[x=vtip, y=Pideal] {results/fig/a_tipspeed.dat};
  \addlegendentry{ideal, $T\vind$}
  \addplot[only marks, mark=*, mark size=2.2pt, cD, forget plot]
      coordinates {(\aVstall,\aPstall)};
  \addplot[only marks, mark=*, mark size=2.2pt, cA, forget plot]
      coordinates {(\aVdesign,\aPdesign)};
  \node[lbl, cD, anchor=south west, fill=white, inner sep=1.2pt] (st)
      at (axis cs:18.6,6.5) {stall boundary};
  \draw[cD, line width=0.4pt] (st.north east) -- (axis cs:\aVstall,{\aPstall-0.7});
  \node[lbl, anchor=south west, fill=white, inner sep=1.2pt] (ds)
      at (axis cs:32,31) {design, $\Ct/\sigma = \aCtDesign$};
  \draw[cA, line width=0.4pt] (ds.south west) -- (axis cs:{\aVdesign+0.3},{\aPdesign+0.7});
\end{axis}
\end{tikzpicture}

\caption{One rotor of the design at its hover thrust, with the geometry fixed
and the tip speed varied (engine data). The induced power $\kappa T\vind$ does
not depend on tip speed; the profile power grows as $\Vtip^3\Cdz(\mathrm{Re})$,
slightly slower than $\Vtip^3$ at low speed because $\Cdz$ falls as the
Reynolds number rises, \cref{eq:recorrection}. The shaft power is therefore
least at the lowest tip speed the blade-loading limit admits, on the stall
boundary $\Ct/\sigma = \aCtMax$ at \SI{\aVstall}{\metre\per\second}. The
design runs at $\Ct/\sigma = \aCtDesign$ and \SI{\aVdesign}{\metre\per\second},
and pays for its stall margin with the difference between
\SI{\aPdesign}{\watt} and \SI{\aPstall}{\watt} of shaft power per rotor.}
\label{fig:tipspeed}
\end{figure}

\subsection{Electrical power}

Motor and electronic speed controller losses are lumped into a single
efficiency $\eta$, so the electrical power drawn by the propulsion system is
\begin{equation}
  \Pelec = \frac{\Pshaft}{\eta}
         = \frac{\kappa T\vind + \Pprof}{\eta}
         = \frac{\Pideal}{\FM\,\eta}
         = \frac{T^{3/2}}{\FM\,\eta\,\sqrt{2\rho A}} .
  \label{eq:pelec}
\end{equation}
The final form is the one carried into \cref{sec:closure}; the intermediate
forms are what the engine actually evaluates, so that $\FM$ is computed from
\cref{eq:fm} rather than assumed.

Adding the non-propulsive load $\Paux$ (display, computer, sensors), the total
electrical power in hover is
\begin{equation}
  P_{\mathrm{tot}}(m) = \frac{(mg)^{3/2}}{\FM\,\eta\,\sqrt{2\rho A}} + \Paux .
  \label{eq:ptot}
\end{equation}

\begin{remark}
Unlike a camera drone, this vehicle carries a payload that consumes power
continuously. $\Paux$ is not negligible: at the designs of \cref{sec:results}
it is \SIrange{14}{16}{\watt} against roughly \SI{80}{\watt} of propulsion, and
because it is constant it converts directly into battery mass, appearing in
\cref{sec:closure} as an addition to the fixed load rather than as a term
scaling with $m$.
\end{remark}
